\centerline{\bf \Huge Правильный перебор I}

\begin{flushright}
        \scriptsize{Здесь может быть Ваша цитата}
\end{flushright}

\problem{\textbf{0}} Выпишите все наборы из трёх цифр, каждая из которых равна 1, 2 или 3, если порядок цифр неважен (т.е. наборы 112 и 121
считаются одинаковыми).

\begin{flushright} \textbf{1 балл}\\ 
\end{flushright}


\problem{\textbf{1}} Кузнечик прыгает вдоль прямой вперёд на 90 см
или назад на 40 см. Может ли он менее чем за 10 прыжков удалиться от начальной точки ровно
на 2 м 50 см?

\begin{flushright} \textbf{2 балла}\\ 
\end{flushright}


\problem{\textbf{2}} МГУ был основан в 1755 году. Сколько раз с момента основания МГУ и до сегодняшнего дня
была дата, которая записывалась с помощью всего двух различных цифр? Каждая дата записывается восемью
цифрами в формате ДД.ММ.ГГГГ (например, 16.07.2025).

\begin{flushright} \textbf{2 балла}\\ 
\end{flushright}


\problem{\textbf{3}} Мультфильм показывали целое число минут. Когда посмотрели в программке время начала и конца показа (минуты и часы по 24-часовой шкале), оказалось, что в записи использованы 8 различных цифр. Какое наименьшее время мог идти мультфильм?

Например: 10:23 - 11:54 - 8 цифр, но они не различны; мультфильм длился 1 час 31 минуту

\begin{flushright} \textbf{2 балла}\\ 
\end{flushright}

